\section{Síntesis y ampliación}

Esta clase mostró en una hora, mediante la ruta de tres pasos de Compound AI: 1) capturar y comprimir las trayectorias de búsqueda, 2) usar la Constitución APAC para gobernar las preguntas y respuestas de múltiples turnos, y 3) aprender restricciones complejas mediante DSPy, el ciclo completo desde el modelado hasta la implementación en ingeniería.

Estas tres partes no son independientes entre sí: el modelo de trace se usa para guiar el flujo de preguntas y respuestas del multi-turn agent, el Json estructurado que genera APAC es a su vez la entrada de restricciones del latent solver en DSPy, y el plan de acción final se envía de vuelta al search trace para una verificación de consistencia, formando un anillo de control auditable.

\begin{table}[H]
\centering
\begin{tabular}{lp{5cm}p{6cm}}
\toprule
Tema & Mecanismo clave & Valor didáctico \\
\midrule
Trayectoria de búsqueda & Trace-aware Transformer + DU Forer drop & Mejora la data/parameter efficiency y refuerza la interpretabilidad \\
Interacción de múltiples turnos & APAC Constitution + agent judge/DPO & Las preguntas de múltiples turnos poseen accuracy/proactivity/credibility \\
Restricciones complejas & Latent cost \(C(Y)\) + solver + green descent & Incorpora retroalimentación diferenciable en escenarios no lineales, sin dejar de lado las restricciones de manufactura \\
\bottomrule
\end{tabular}
\caption{Componentes clave y valor del método de tres pasos de Compound AI}
\end{table}

\subsection{Preguntas de investigación y extensiones}

\paragraph{Gestión de versiones de la capa de entrenador de IA}
Cuando el search trace, el agent APAC y el solver de DSPy iteran al mismo tiempo, es fácil que surja una «desincronización entre los parámetros del trace y el juicio del agent». Se recomienda empaquetar juntos el ID del trace, la versión del modelo de agent judge y la solver configuration en un registro de release, y ejecutar una prueba automatizada antes de cada despliegue, para asegurar que la nueva versión del judge todavía pueda reconocer los trace antiguos y evitar una unstable inference.

\paragraph{Generalización entre dominios y calibración del judge}
El agent judge actual destaca en tareas de travel intent, pero en ámbitos de alto riesgo como la medicina o las finanzas su comprensión del goal podría ser completamente distinta. Una estrategia de extensión viable es entrenar un mixture-of-experts para el judge o introducir un domain embedding que le permita adaptarse a un new intent, y a la vez alinear el judge reward con la human satisfaction signal (NPS, tasa de clientes recurrentes), en lugar de depender únicamente de la plan optimality.

\paragraph{Fusión multimodal y materiales de apoyo}
El expositor mencionó que «los slides, notebooks y photos» son herramientas poderosas para reforzar la explicación. Idealmente, el search trace debería poder referenciar las tablas/fórmulas de los slides (como el sequencing graph de Soang); cuando aparezca un diagram clave en el video, capturar automáticamente el fotograma e insertarlo en la note, manteniendo una doble evidencia visual y textual.

\begin{itemize}
    \item ¿Cómo comprimir aún más la cantidad de tokens del trace para reducir el costo de inferencia, manteniendo al mismo tiempo la interpretabilidad de múltiples turnos?
    \item ¿Cómo puede transferirse la capacidad de generalización del agent judge a ámbitos de alto riesgo como la medicina o las finanzas, y no solo al travel intent?
    \item ¿Pueden las restricciones de manufactura diferenciadas abstraerse automáticamente en una latent representation, de modo que se transfieran sin fricciones a nuevas industrias (por ejemplo, el diseño de chips)?
\end{itemize}

Además de las preguntas teóricas, este Compound AI también debe consolidarse de manera continua en production: mantener el audit log de trace/plan, alinear el DPO reward con la satisfacción del usuario, e iterar el modelo de judge en distintas region, para así asegurar la robustez del sistema a largo plazo.

\subsection{Lecturas adicionales}
\begin{itemize}
    \item «Beyond A*: Better Planning with Transformers via Search Dynamics Bootstrapping», útil para entender cómo las trayectorias de búsqueda mejoran el desempeño de la planificación.
    \item «Dualformer: Controllable Fast and Slow Thinking by Learning with Randomized Reasoning Traces», que corresponde al trace drop y a la velocidad de razonamiento controlable.
    \item «Composing Global Optimizers to Reasoning Tasks via Algebraic Objects in Neural Nets», que analiza la optimización global composicional en redes neuronales.
    \item «SurCo: Learning Linear Surrogates For Combinatorial Nonlinear Optimization Problems», que muestra una ruta de aprendizaje de surrogate para problemas de optimización combinatoria no lineal.
    \item APAC Constitution (2024), libro blanco interno que define en detalle los métodos de evaluación de Accuracy/Proactivity/Adaptivity/Credibility.
\end{itemize}

\subsection{Resumen de la sección}
El método de tres pasos de Compound AI no es un experimento aislado, sino una combinación orientada a production: fusiona los visual slides, los time-stamped frames y los trace logs en un mismo pipeline, para que la audiencia no solo entienda los principios del método, sino que también pueda dar seguimiento a la implementación concreta en ingeniería.

\end{document}
